\section{Normalized adic characters and their actual coefficient rings}
\label{sec:adic-realization}

This section closes the one-dimensional adic-character supplier left
explicit in AR. The input is a specified global reciprocity map
\[
 \operatorname{Art}:\mathbb A_E^\times\longrightarrow G_E^{\mathrm{ab}},
 \qquad
 \ker\operatorname{Art}
   =\overline{E^\times E_{\infty,+}^\times},                 \tag{GA.1}
\]
with the open quotient, inertia, decomposition and
\emph{inverse arithmetic Frobenius} clauses of RC3. Global class field
theory itself is retained. All conclusions below use this same map.
The algebraic closure of $\mathbb Q_\ell$, denoted
$\overline{\mathbb Q}_\ell$, carries its extended $\ell$-adic norm;
it is not silently replaced by its completion. The field isomorphism
$\iota:\overline{\mathbb Q}_\ell\simeq\mathbb C$ is algebraic and is
never assumed continuous.

Our coefficient datum consists of a DVR $O$, finite as a
$\mathbb Z_\ell$-module, its actual maximal-ideal-adic topology, and
an embedding
$\kappa:\operatorname{Frac}(O)\hookrightarrow\overline{\mathbb Q}_\ell$
that agrees with the standard map on $\mathbb Z_\ell$.
Write $g=\kappa|_O$. A Hecke character is a continuous
$\chi:\mathbb A_E^\times\to\mathbb C^\times$ trivial on $E^\times$.
For $n:\operatorname{Hom}(E,\mathbb C)\to\mathbb Z$,
\texttt{IsAlgPos} means
$\chi_\infty(a)=\prod_\tau\tau(a)^{n_\tau}$ for totally positive
$a\in E^\times$. A normalized adic realization is a continuous
$r:G_E\to O^\times$ satisfying
\[
 \kappa(r(\sigma))
 =\iota^{-1}\!\left(\chi(x)F_{\infty,n}(x)^{-1}\right)
       F_{\ell,n}(x)
 \quad\text{if }\operatorname{Art}(x)=\bar\sigma.            \tag{GA.2}
\]
Here the two finite products $F_{\infty,n}$ and $F_{\ell,n}$ are
defined in GA4. This is the literal \texttt{IsGal} identity,
including its coefficient embeddings.

\paragraph{GA1: the norm detects the actual DVR and its units.}
For the coefficient datum above, every $a\in O$ satisfies
$|g(a)|_\ell\leq1$, every unit has norm $1$, and
\[
 x\in O\ \Longleftrightarrow\ |\kappa(x)|_\ell\leq1,
 \qquad
 x\in O^\times\ \Longleftrightarrow\ |\kappa(x)|_\ell=1.      \tag{GA.3}
\]
We identify $O$ with its image in its fraction field in (GA.3).
Moreover $\ell\in\mathfrak m_O$.

Finiteness over $\mathbb Z_\ell$ makes $a$ integral. In a monic
equation for $a$, if $|g(a)|_\ell>1$, the leading term has larger
norm than every lower term, since their coefficients have norm at
most $1$. The ultrametric inequality rules out their sum being zero.
Applying this to a unit and its inverse gives norm $1$.
The element $\ell$ has norm $\ell^{-1}<1$ and is nonzero, so it is
a nonunit. Choose a uniformizer $\pi$ and write
$\ell=u\pi^e$, where $u\in O^\times$ and $e\geq1$.
Then
$0<|g(\pi)|_\ell=\ell^{-1/e}<1$.
Every nonzero fraction is uniquely a unit times an integral power of
$\pi$. Its norm is therefore $|g(\pi)|_\ell^a$, and it belongs to
$O$, respectively $O^\times$, exactly when $a\geq0$,
respectively $a=0$. This proves both directions, rather than just
boundedness of integral elements.

\paragraph{GA2: the coefficient fraction field has closed image.}
The image $L=\kappa(\operatorname{Frac}(O))$ is a finite-dimensional
$\mathbb Q_\ell$-subspace of $\overline{\mathbb Q}_\ell$, hence is
closed in that normed field.
Choose $\mathbb Z_\ell$-module generators $s_1,\ldots,s_d$ of $O$.
Every nonzero denominator $b=u_b\pi^a$ divides a power of $\ell$:
with $\ell=u\pi^e$, $a\geq0$,
\[
 b\bigl(u_b^{-1}u^a\pi^{(e-1)a}\bigr)=\ell^a.
\]
Consequently $\kappa(a_0/b)$ is a $\mathbb Q_\ell$-linear combination
of the $g(s_j)$. Conversely every such combination belongs to $L$:
a $\mathbb Q_\ell$ scalar is a fraction of elements of
$\mathbb Z_\ell$, and $\kappa$ respects their maps. Thus
$L=\sum_j\mathbb Q_\ell g(s_j)$.

For completeness, finite-dimensional subspaces over
$\mathbb Q_\ell$ are closed even when the ambient normed space is
incomplete. On a basis of such a subspace the coordinate sup norm
and the induced norm are equivalent. The lower estimate follows
by taking the minimum of the induced norm on the compact coordinate
unit sphere, where linear independence prevents zero; the upper
estimate follows from the triangle inequality. Cauchy sequences
therefore have convergent coordinates in complete $\mathbb Q_\ell$.
Their limits stay in the subspace, which proves closedness.
This is the closed-image fact needed for the quotient argument in
GA8; a merely algebraic subfield of
$\overline{\mathbb Q}_\ell$ would not suffice.

\paragraph{GA3: the embedding reflects topology and coefficient maps are continuous.}
The injection $g:O\to\overline{\mathbb Q}_\ell$ is a topological
embedding for the specified adic topology. Put
$c=|g(\pi)|_\ell\in(0,1)$. The inclusions
\[
 g(\pi^aO)\subseteq\{z:|z|_\ell\leq c^a\},
 \qquad
 |g(x)|_\ell<c^a\Longrightarrow x\in\pi^aO                \tag{GA.4}
\]
compare neighborhood bases. For the second, apply GA1 to
$x/\pi^a$. Since $c^a\to0$, they prove the assertion.
Continuity of an $O$-valued map can thus be tested after $g$.
For a homomorphism from a topological group into $O^\times$,
continuity of its underlying $O$-valued map also gives continuity
in the unit topology: the inverse coordinate is its value at
the inverse group element.

More generally, let $O'$ be an adic topological local ring and
$t:O\to O'$ a unital ring homomorphism, with
$t(\ell)\in\mathfrak m_{O'}$. This map is continuous.
Indeed
$t(\pi)^{ea}=t(u)^{-a}t(\ell)^a\in\mathfrak m_{O'}^a$;
the image of $\pi^{ea}O$ lies in $\mathfrak m_{O'}^a$.
This uses the actual adic topologies, and does not assume $t$
is a local map beforehand. In particular it applies to the
coefficient extensions used in GA19 and GA21.

\paragraph{GA4: the archimedean and adic algebraic products.}
For an idele $x$ define
\[
 F_{\infty,n}(x)=
 \prod_{v\mid\infty}\ \prod_{\substack{j:E_v\to\mathbb C\\j\ {\rm continuous}}}
       j(x_v)^{\,n_{j|_E}},
 \qquad
 F_{\ell,n}(x)=
 \prod_{v\mid\ell}\ \prod_{\substack{j:E_v\to\overline{\mathbb Q}_\ell\\
                                  j\ {\rm continuous}}}
       j(x_v)^{\,n_{\iota j|_E}}.                         \tag{GA.5}
\]
They are continuous multiplicative maps into the corresponding
unit groups. Each indexing set is finite: restriction to the dense
subfield $E$ is injective, and a number field has finitely many
embeddings. There are no continuous embeddings from $E_v$ into
$\overline{\mathbb Q}_\ell$ when $v\nmid\ell$.
For otherwise choose an integer divisible by the residue
characteristic of $v$ and prime to $\ell$; its powers tend to zero
in $E_v$ while their images have $\ell$-adic norm $1$.

Restriction identifies the union of the archimedean embedding sets
with all embeddings $E\to\mathbb C$. Likewise the union of the
$\ell$-adic embedding sets identifies, after $\iota$, with those
same embeddings. To see the latter without assuming continuity of
$\iota$, begin with $\tau:E\to\overline{\mathbb Q}_\ell$.
The valuation of $\tau$ gives a place $v\mid\ell$; equivalence of
the discrete valuations shows continuity on $E$ for $v$.
The image lies in the finite extension
$\mathbb Q_\ell(\tau(E))$, which is complete by the argument of
GA2, so $\tau$ extends uniquely to $E_v$. Conversely a continuous
extension determines its place from the induced valuation, and
density gives uniqueness. At infinity the same assertion is
the extension to $\mathbb R$ or $\mathbb C$, with complex
conjugation giving the second embedding of a complex completion.
Consequently, for $a\in E^\times$,
\[
 F_{\infty,n}(a)=\prod_\tau\tau(a)^{n_\tau},\qquad
 F_{\ell,n}(a)=\iota^{-1}\!\left(\prod_\tau\tau(a)^{n_\tau}\right).
                                                                  \tag{GA.6}
\]
The adic product is $1$ on purely infinite ideles; the infinite
product is $1$ on purely finite local ideles.

\paragraph{GA5: removing the infinite type gives an open kernel.}
Suppose $\chi$ satisfies \texttt{IsAlgPos} for $n$ and set
$\theta=\chi F_{\infty,n}^{-1}$. Then $\theta$ is continuous and
trivial on $E_{\infty,+}^\times$, and its kernel in the idele group
is open.

First, the connected component of $1$ in
$E_\infty^\times$ is the product of positive real factors and
complex multiplicative factors; it is open. Weak approximation
makes $E^\times$ dense in $E_\infty^\times$. Its intersection with
that component consists of totally positive elements and is dense
there. The defining \texttt{IsAlgPos} equality and continuity
therefore give $\theta=1$ on the whole component.

Every continuous homomorphism from the finite idele group to
$\mathbb C^\times$ has open kernel. Here is the needed elementary
argument. There is a neighborhood of $1$ on which all image values
lie in $|z-1|<1/2$. Inside it choose an open subgroup of the form
\[
 \{u:u-1\in N\widehat{\mathcal O}_E,\quad
           u^{-1}-1\in N\widehat{\mathcal O}_E\},\quad N\geq1.
                                                                  \tag{GA.7}
\]
Such subgroups form a neighborhood basis by the restricted-product
topology. If $z$ and all its powers satisfy $|z^a-1|<1/2$,
then $z=1$: the geometric sum $S_a=\sum_{h<a}z^h$ satisfies
$|S_a-a|\leq a/2$, so $|S_a|\geq a/2$, while
$|(z-1)S_a|=|z^a-1|<1/2$, a contradiction for large $a$ if
$z\ne1$. All powers of an element of (GA.7) stay in that
subgroup, proving it is killed. Combining this finite open
subgroup with $E_{\infty,+}^\times$ proves the claim.
It follows that $\iota^{-1}\theta$ is continuous: an arbitrary
homomorphism with open kernel is locally constant, irrespective
of the topology of $\iota$.

\paragraph{GA6: construction and uniqueness of the normalized adic character.}
Let $\chi,n$ satisfy \texttt{IsAlgPos}. Define a continuous
idele character
\[
 \Psi(x)=\iota^{-1}\!\left(\chi(x)F_{\infty,n}(x)^{-1}\right)
                          F_{\ell,n}(x).                 \tag{GA.8}
\]
Continuity follows from GA4 and GA5. On a principal idele,
$\chi(a)=1$ and (GA.6) cancels the other two factors.
On $E_{\infty,+}^\times$ the first factor is $1$ by GA5 and
the second is $1$ by GA4. The closed kernel of $\Psi$
therefore contains the kernel (GA.1). It descends uniquely to
a continuous character of $G_E^{\mathrm{ab}}$, since
$\operatorname{Art}$ is an open quotient map. Composing with
$G_E\to G_E^{\mathrm{ab}}$ gives a continuous
$r_\chi:G_E\to\overline{\mathbb Q}_\ell^\times$ satisfying
(GA.2).

If another character satisfies (GA.2), choose an idele mapping
to each $\bar\sigma$ by surjectivity of $\operatorname{Art}$.
The displayed equality gives the same value on $\sigma$.
This proves uniqueness. Thus the source's published
\texttt{HeckeLAdicCharacter} statement follows from the retained
\texttt{GlobalArtinReciprocity}; no second character-existence
hypothesis is required.

\paragraph{GA7: local reciprocity, inertia, and the Frobenius inverse.}
If $\operatorname{Art}(x_v)=\bar\sigma$, then
\[
 r_\chi(\sigma)=\iota^{-1}(\chi_v(x_v))
       \prod_jj(x_v)^{n_{\iota j|_E}}.                    \tag{GA.9}
\]
This is (GA.2) for a local finite idele. The product is $1$
when $v\nmid\ell$, or when $n=0$. Under either condition,
the inertia-image clause in (GA.1) gives exactly
\[
 r_\chi(I_v)=\iota^{-1}(\chi_v(\mathcal O_v^\times)).
                                                                  \tag{GA.10}
\]
In particular $r_\chi$ is unramified at $v$ if and only if
$\chi$ is, and the two inertia images have the same exponent
and, when finite, the same order. The decomposition clause,
continuity and closure similarly show that $r_\chi$ is
trivial on the whole decomposition group if and only if
$\chi_v$ is the trivial local character.

For arithmetic Frobenius $\sigma$ at any valuation above an
unramified $v$, the specified Artin convention reads a
uniformizer as $\bar\sigma^{-1}$. Consequently
\[
 \iota(r_\chi(\sigma))=\chi(v)^{-1}
       \quad(v\nmid\ell\text{ or }n=0).                 \tag{GA.11}
\]
These assertions also hold for $O$-valued realizations after
applying $g$. Multiplication and inversion of normalized
characters add and negate $n$, because the two products in
(GA.5) have exactly those laws. No inverse is added to the
local unit formula (GA.9).

\paragraph{GA8: descent into the specified coefficient DVR.}
Suppose a field $K$ has embeddings
$\epsilon:K\hookrightarrow\mathbb C$ and
$f:K\hookrightarrow\operatorname{Frac}(O)$ satisfying
$\iota\kappa f=\epsilon$, and that $\chi(v)\in\epsilon(K)$
for all but finitely many finite places. Then the character
of GA6 takes its values in $g(O^\times)$ and gives the exact
\texttt{exists\_isGal} conclusion.

Let $L$ be the closed field image from GA2. The subgroup
$L^\times\subseteq\overline{\mathbb Q}_\ell^\times$ is
closed, since it is the inverse image of $L$ under the
continuous underlying-value map. The quotient by it is a
Hausdorff commutative topological group. Compose $r_\chi$
with this quotient. Outside the finitely many excluded,
ramified and $\ell$-adic places, (GA.11) belongs to
$\kappa f(K^\times)$, so the quotient character kills every
Frobenius. RC4 proves that the quotient character is
identically $1$. Thus every value of $r_\chi$ belongs to $L$.

The absolute Galois group is compact. A continuous character
from a compact group into a normed field has values of norm
$1$: its image is bounded, and a value with norm greater
than $1$ would have unbounded positive powers; apply the
same argument to its inverse. By GA1 the unique preimage
of each value lies in $O^\times$. Uniqueness makes these
preimages multiplicative. GA3 makes the resulting $O$-valued
map continuous, and the unit inverse coordinate is continuous
as well. Its embedded character is precisely $r_\chi$, so it
satisfies (GA.2) for the original $\chi,n,\kappa,\iota$.
This constructs the normalized character itself; it does not
replace it by a finite-order character.

\paragraph{GA9: finite-order characters and exact recovery.}
If $\chi^m=1$ with $m>0$, then it satisfies
\texttt{IsAlgPos} with $n=0$. Indeed every totally positive
infinite idele has an $m$th root in the product of positive
real and nonzero complex factors; evaluating $\chi$ on
that power gives $1$. Formula (GA.2) then gives
$r_\chi^m=1$. Its image is a finite set of roots of unity,
so its kernel is open: the complement of $1$ in that
finite set is closed, and continuity pulls its complement
back to the kernel.

Conversely, let $r:G_E\to\overline{\mathbb Q}_\ell^\times$
be continuous with $r^m=1$. It descends through topological
abelianization, since its target is commutative and its
kernel is closed. Set
$\chi(x)=\iota(\bar r(\operatorname{Art}(x)))$.
This is continuous even though $\iota$ need not be:
$r$ has open kernel and the resulting character is locally
constant. It kills principal ideles, has order dividing
$m$, and satisfies (GA.2) with $n=0$.
Surjectivity onto the abelianization proves uniqueness.
The local identities of GA7 hold at \emph{all} finite places
in this case, including those above $\ell$.

\paragraph{GA10: anticyclotomic symmetry follows from the actual character.}
Suppose $E/F$ is the host quadratic CM extension with conjugation
$c$, and $\widetilde c$ is the supplied automorphism of the
algebraic closure extending it. Retain the host's finite-place
unramified equality
$v_P(a)=v_{P\cap F}(a)$ for $a\in F^\times$, and let
$\chi$ be trivial on the ideles of $F$. Every normalized
realization satisfies
\[
 r_\chi(\widetilde c\sigma\widetilde c^{-1})
      =r_\chi(\sigma)^{-1}.                             \tag{GA.12}
\]
The map $\sigma\mapsto\widetilde c\sigma\widetilde c^{-1}$
is continuous for the Krull topology: fixing the conjugates
of a finite field basis forces the conjugate automorphism
to fix the original finite field.

At an unramified pair $P,\bar P$ outside the $\ell$-adic
places, restriction of $\chi$ to a uniformizer of the
underlying place of $F$ gives
$\chi(P)\chi(\bar P)=1$. If $P=\bar P$, the same restriction
and the retained valuation equality give $\chi(P)=1$,
which is enough for this product identity. If they are
different, the two primes are all primes over that place,
each with exponent one. Conjugating a valuation and its
arithmetic Frobenius preserves the residue-field cardinality
and gives Frobenius at $\bar P$. Equation (GA.11) therefore
makes the continuous character
$\sigma\mapsto r_\chi(\widetilde c\sigma\widetilde c^{-1})
r_\chi(\sigma)$ trivial on every Frobenius outside a finite
conjugation-stable set. RC4 proves (GA.12). This proves the
source's \texttt{AntiHom} property with its supplied
conjugation, rather than choosing a new involution.

\paragraph{GA11: the CM parity calculation gives integral infinite types.}
Use the host CM type $\Phi=\{\phi_\tau\}$ of AH and HF.
Every complex embedding of $E$ occurs exactly once among
$\phi_\tau,\bar\phi_\tau$: equality of their infinite places
first gives equality of $\tau$, and total complexity excludes
an embedding fixed by conjugation. For $z\ne0$ and $a\in\mathbb Z$,
\[
 (z/|z|)^{2a}=z^a\bar z^{-a}.                            \tag{GA.13}
\]
Thus a character with argument exponent $2a_\tau$ has
algebraic exponents $(a_\tau,-a_\tau)$ on this pair.
In particular the HF character $\mu_{0,j}$, whose argument
exponent is $-2e_{j,\tau}$, is algebraic when the prescribed
$e_{j,\tau}$ are integers.

For an odd integer $r=2h+1$ and the HF character $\eta$ with
argument exponent $1$, the infinite contribution of
$\eta^{-1}|\cdot|^{-r/2}$ at a complex place is
\[
 (z/|z|)^{-1}|z|^{-r}
       =z^{-h-1}\bar z^{-h}.                            \tag{GA.14}
\]
Here the local absolute-value factor at a complex place
counts both embeddings. Hence this character is algebraic
with integral exponents $(-h-1,-h)$.
The source has $r=k+m-1$, with $k,m$ even, so $r$ is odd;
$t$ is separately odd. The same argument applies to $t$.
It does not identify the global chain length $m$ with any
local polynomial degree.

\paragraph{GA12: odd normalization preserves the specified field of values.}
Write
\[
 B_r(\mu,\eta)=\mu\,\eta^{-1}|\cdot|^{-r/2}.
 \qquad
 B_r(\mu,\eta)(P)=\mu(P)\eta(P)^{-1}N(P)^{r/2}.           \tag{GA.15}
\]
The second equality uses $|\varpi_P|=N(P)^{-1}$.
Suppose $\mu(P)\in\epsilon(K)$ cofinitely often, and that
$\eta(P)N(P)^{-1/2}\in\epsilon(K)$ at every unramified
place, the literal \texttt{ValuesIn} normalization for
$\eta$. If $r=2h+1$, choose $y,z\in K^\times$ with
$\epsilon(y)=\mu(P)$ and
$\epsilon(z)=\eta(P)N(P)^{-1/2}$. Then
\[
 B_r(\mu,\eta)(P)=\epsilon(yz^{-1}N(P)^h).                \tag{GA.16}
\]
The absolute norm is a positive integer, so its integral
power belongs to the prime field. This gives
\texttt{AlmostIn} for the \emph{same} field $K$.
Products and inverses preserve this cofinite condition
by intersecting the finite exceptional sets.
On a local unit, if $\eta$ is unramified,
\[
 B_r(\mu,\eta)_P(u)=\mu_P(u).                            \tag{GA.17}
\]
The norm factor and $\eta_P(u)$ are both $1$.

\paragraph{GA13: construction of the exact shift and block characters.}
Fix the actual $D$, its host character frame $C$, an index
$i$ of the supplied $\mathrm{NTw}$, a side $b$, and a
\texttt{Reading} $(\kappa,\iota)$ for its ring $O_0(i)$.
Thus
\[
 \iota\kappa\,\operatorname{embSide}_{i,b}(x)
       =\operatorname{emb}_{LL_0}(x)
 \quad(x\in LL_0),                                     \tag{GA.18}
\]
and $\kappa$ has the specified $\mathbb Z_\ell$ compatibility.
Retain \texttt{ValLL0} for this actual frame; HF16 supplies
it when its value field is contained in this actual $LL_0$.
There exist normalized $O_0(i)^\times$-valued characters
for
\[
 \operatorname{shiftChar}(C)=B_t(1,\eta),\qquad
 \operatorname{blockChar}(C,\ell,j)
        =B_r(\mu_{0,j}\delta_\ell,\eta),                 \tag{GA.19}
\]
where $\delta_\ell=\delta'$ if $\ell=p_\delta$, and
$\delta_\ell=\delta$ otherwise.

GA11 gives the integral infinite type of the shift.
The auxiliary character is finite order, hence algebraic
of type zero by GA9; GA11 then gives the block's type.
The odd $t$ and odd $r$ calculations of GA12 put their
values cofinitely in the actual $LL_0$.
Apply GA8 with $K=LL_0$, $f=\operatorname{embSide}_{i,b}$
and (GA.18). This proves the literal
\texttt{Triple.exists\_shift} and
\texttt{Triple.exists\_blockChar} endpoints. The character
and its integral type are chosen after the actual reading;
the construction does not select a coefficient place or
produce a twist index.

\paragraph{GA14: changing a reading does not change the integral character.}
Let $\rho,\rho'$ be two readings of the same
$(i,b,O_0(i),\operatorname{embSide}_{i,b})$. If
$y\in O_0(i)$ and
$\rho_{\mathbb C}(y)=\operatorname{emb}_{LL_0}(x)$, then
injectivity of $\rho.\iota$ and $\rho.\kappa$ gives the
fraction-field equality
$y=\operatorname{embSide}_{i,b}(x)$.
Applying $\rho'$ gives the same complex value.

Suppose an integral character $r$ realizes the block for
$\rho$. Construct $r'$ for $\rho'$ by GA13. At cofinally
many unramified places the inverse block value lies in
the common $LL_0$ by GA12. Equation (GA.11), the preceding
fraction-field equality and injectivity show
$r(\mathrm{Frob})=r'(\mathrm{Frob})$.
Both are continuous in the same adic ring; RC4 gives
$r=r'$. It follows that $r$ also realizes the block for
$\rho'$, with its appropriate infinite type.
The tame-generator equality in \texttt{IsChars} is
independent of the reading too: lift the specified
inertia generator to a local unit using the same Artin
map and apply (GA.9) to both readings. This proves
\texttt{isGal\_of\_reading} and \texttt{isChars\_of\_reading}
without treating arbitrary complex embeddings as continuous.

\paragraph{GA15: identification with the already installed first character.}
Retain the literal \texttt{NormalizedCharacterCompatibility}
record for $D$ and its $NChiData$, and the three identities
identifying its $\eta$, $\mu_0(j+1)$ and $\delta_i$ with
the same frame $C$. For every existing $i,b,j$ and every
reading $\rho$, the \emph{already installed}
$D.\mathrm{ntw}.\chi_1(i,j+1,b)$ realizes the exact block
in (GA.19).

Indeed GA13 constructs a comparison character $r$.
The retained \texttt{chi1} reading says, away from a finite
set, that for some $x_P\in LL_0$ and unramified Hecke value
$\alpha_P$,
\[
 \operatorname{emb}_{LL_0}(x_P)
       =\alpha_P^{-1}\sqrt{N(P)}^{-r},\qquad
 \chi_1(\sigma)=\operatorname{embSide}_{i,b}(x_P)           \tag{GA.20}
\]
in the fraction field at every Frobenius above the specified
valuation. The three character identities make
$\alpha_P=(\mu_{0,j}\delta_\ell\eta^{-1})(P)$.
Since $N(P)>0$, (GA.15) shows that the first expression
in (GA.20) is exactly the inverse block value.
Equations (GA.11) and (GA.18) therefore identify the two
embedded Frobenius values.

The source initially specifies one valuation above $P$.
Any other is its Galois conjugate; conjugating the
Frobenius equation gives Frobenius at the specified one,
and homomorphisms into commutative groups are unchanged
by this conjugation. Thus equality holds at all valuations
required by RC4. The installed character is continuous
by its $\mathrm{NTw}$ field, so RC4 and injectivity give
$r=\chi_1$. This proves
\texttt{firstTwistCharacter\_eq\_ladic\_heckeCharacter}.
The proof consumes the normalized reading record; existence
of a newly constructed $r$ alone would not identify an
arbitrarily preinstalled $\chi_1$.

\paragraph{GA16: exact auxiliary inertia and both side families.}
Let the actual HF frame and the integral $e_{j,\tau}$ be
fixed. Assume the two type primes split completely as in
the source, and retain both coefficient embeddings
$f,f':K\to\operatorname{Frac}(O)$ with
\[
 \iota\kappa f=\epsilon,\qquad
 \iota\kappa f'=\overline{\epsilon}.
\]
Put $\iota_+=\iota$, $\iota_-=\overline{\iota}$.
Then $\iota_b\kappa f_b=\epsilon$ for $f_+=f,f_-=f'$.
Choose the auxiliary prime
$p=p_\delta'$ when $\ell=p_\delta$, and $p=p_\delta$
otherwise. It is prime, distinct from $\ell$ and both
type primes, and $p\equiv1\pmod{15}$.
GA8, GA9 and GA13 construct $\Delta$ for $\delta_\ell$
and $\chi_{j,b}$ for each block with reading $\iota_b$.

The kernel of $\Delta$ is open, and GA10 makes it
anticyclotomic. GA7 and the HF12 exact local exponent
give
\[
 \Delta\text{ ramifies exactly above }p,\qquad
 \#\Delta(I_P)=15\quad(P\mid p).                        \tag{GA.21}
\]
For the order assertion the inertia image is a subgroup
of the $15$th roots in a field, hence finite cyclic.
Its exponent is $15$ because the local exponents divide
$a$ exactly when $15\mid a$; a finite cyclic group's
exponent equals its order. Also $N(P)=p^{f_P}$, so
$N(P)\equiv1\pmod{15}$.

At $P\mid p$, (GA.17) kills the $\mu_0$ factor and
leaves $\delta_\ell$. Since finite-order complex values
have conjugate equal to inverse, (GA.9) gives
\[
 \chi_{j,+}|_{I_P}=\Delta|_{I_P},\qquad
 \chi_{j,-}|_{I_P}=\Delta^{-1}|_{I_P}.                   \tag{GA.22}
\]
Away from $\ell$, the two type primes and $p$, these
characters are unramified. At a type prime $g_k\ne\ell$,
their inertia values have order dividing $g_k-1$, by
tameness and complete splitting, as detailed in GA17.
To define the literal family for every natural $j$,
use the genuine finite family at
$\min(j-1,m-1)$; $m\geq2$ makes this index legitimate,
and it recovers each prescribed $j+1$ for $j\in\operatorname{Fin}(m)$.
This proves every clause of
\texttt{HeckeTwistCharactersDeltaChi}.

\paragraph{GA17: tame and wild local unit orders.}
At a finite place with residue cardinality $q$, reduction
gives $u^{q-1}\equiv1$ for every local unit $u$.
A tame Hecke character therefore satisfies
$\chi_P(u)^{q-1}=1$.
Conversely, suppose a continuous local character kills
all $(q-1)$st roots of unity in the local field.
Its restriction to local units has open kernel: choose
a sufficiently small congruence subgroup mapped into
$|z-1|<1/2$ and apply the power argument of GA5.
Hensel's lemma applied to $X^{q-1}-1$, whose derivative
is a unit, gives a root $\zeta$ with the same residue as
$u$. Then $v=u/\zeta$ belongs to $1+\mathfrak m_P$.
Writing $q$ in the maximal ideal and using
\[
 v^q-1=(v-1)\sum_{a<q}v^a,\qquad
 \sum_{a<q}v^a=q+\sum_{a<q}(v^a-1)
 \in\mathfrak m_P,
\]
shows that each $q$th power raises the additive depth
by at least one. Hence $v^{q^a}$ eventually lies in
the open kernel of the character, and
$\chi_P(u)^{q^a}=1$ for some $a$.
This proves the two exact \texttt{orders\_on\_units}
statements, including dyadic places.

We will also need the Galois version. Let a continuous
commutative character have inertia image of exponent
$m$ prime to the residue characteristic, and let $\gamma$
be a specified tame generator. Its image kills wild
inertia, by the finite ramification filtration argument
of AR3. Continuity and the finiteness of the $m$th roots
allow this argument in a finite Galois quotient: choose
an open normal subgroup whose image avoids every
nonidentity $m$th root. By the definition of a tame
generator, inertia is the closure of the group generated
by $\gamma$ and wild inertia. The preimage of the finite
cyclic group generated by the character value at $\gamma$
is closed and contains those generators. Thus every
inertial value is a power of that value. This proves
\texttt{mem\_zpowers\_of\_tameGen} for the continuous
characters used below.

\paragraph{GA18: a uniform calibration of the specified tame generator.}
Fix a type prime $p=g_k$, a place $P$ in its chosen half,
a valuation above it and its specified tame generator
$\gamma_P$. Complete splitting gives residue field
$\mathbb F_p$. Lift $\gamma_P$ by the Artin inertia
clause to a local unit $u$, and denote its residue by
$a\in\mathbb F_p^\times$. The HF local law, first on
rational residue lifts and then on all units by tameness,
reads
\[
 \mu_{0,j,P}(u)=\omega_k(a)^{-2(j+1)}.                  \tag{GA.23}
\]
We claim that $a^2$ generates the subgroup of squares.
Use the row $j=0$ and an adic realization of $\mu_{0,0}$
at the distinct auxiliary prime $p_\delta$; GA6 supplies
it, without any coefficient-field restriction. Its
inertia image at $P$ has exponent dividing $p-1$.
For every $b\in\mathbb F_p^\times$, lift $b$ to a rational
local unit and then to inertia. GA17 makes its character
value a power of the value at $\gamma_P$. Equation
(GA.23) and injectivity of $\omega_k$ imply
$b^2\in\langle a^2\rangle$.

In a cyclic group of order $p-1$, this implies there is a
generator $a_0$ with $a_0^2=a^2$. To check the possible
factor of two, write the group as $\langle z\rangle$,
$a=z^d$. The condition is
$\gcd(2d,p-1)=\gcd(2,p-1)$.
If $p-1$ is odd, $a$ already generates.
If it is even and $\gcd(d,p-1)=1$, the same is true.
Otherwise $\gcd(d,p-1)=2$ and $(p-1)/2$ is odd;
$a z^{(p-1)/2}$ generates and has the same square.

The two injective characters $\omega_k$ and the specified
Teichmuller map of $\mathbb F_p^\times$ have the same
image, namely all $(p-1)$st roots of unity. Transport
$a_0$ through these injections, obtaining a generator
$a_P$ read by the specified Teichmuller map. For
\emph{every} $\ell\ne p$, every allowed coefficient ring
and reading, every $j$, and every $\xi$ unramified at $P$,
any realization of $B_r(\mu_{0,j}\xi,\eta)$ satisfies
\[
 \iota\kappa(\chi(\gamma_P))
   =\operatorname{Teich}\bigl(a_P^{-2(j+1)}\bigr).        \tag{GA.24}
\]
This follows directly from (GA.9), (GA.17), (GA.23)
and $a_0^2=a^2$. Crucially $a_P$ is chosen before all
these coefficient data and $j$. This proves
\texttt{TameInertiaGeneratorCharacterSpec}; setting
$c(k,P,j)=a_P^{2(j+1)}$ gives the actual
\texttt{exists\_tameTie}, for a new such array.
It does not identify that array with an arbitrary
preinstalled $D.\mathrm{nd}.\chi_T$.

\paragraph{GA19: the toric quotient character and its integral Frobenius values.}
Let $\vartheta$ be a Hecke character trivial on $F$,
with even integral argument exponents. Assume the
specified $K$ contains the values of $\mu_{0,j}$,
$\delta_\ell$ and $\vartheta$.
GA8 and GA11 construct the normalized realization
$\lambda_0$ of
$\mu_{0,j}\delta_\ell\vartheta^{-1}$, using the selected
side reading. For any coefficient map $O\to O'$ satisfying
GA3, let $\lambda$ be its image in $O'^\times$.
It is continuous and anticyclotomic by GA10.
If $S_0$ contains the $\ell$-adic, type and auxiliary
places and all ramification of $\vartheta$, then
$\lambda$ is unramified outside $S_0$ by GA7.

Suppose at the auxiliary places $\vartheta_P$ kills
the $(N(P)-1)$st roots of unity. On local units the
quotient of $\lambda$ by the mapped
$\Delta^{\epsilon_b}$ is read as $\vartheta_P^{-1}$,
where $\epsilon_+=1,\epsilon_-=-1$.
GA17 therefore proves, for every $\sigma\in I_P$,
\[
 \bigl(\lambda(\sigma)\,
          \Delta'(\sigma)^{-\epsilon_b}\bigr)^{N(P)^a}
        =1\quad\text{for some }a\geq0.                 \tag{GA.25}
\]
This is the exact power-of-residue-cardinality conclusion,
not a claim that the quotient is trivial.
Define $X(P)$ as the image of $\lambda_0$ at the source's
chosen Frobenius at $P$. If $P\notin S_0$, equation
(GA.11) shows the value is independent of every choice
of valuation and Frobenius. Thus $X(P)$ lies in $O'$,
with the literal equality to $\lambda(\sigma)$ for
all such choices. These constructions prove
\texttt{ToricHeckeQuotientCharacter}.

\paragraph{GA20: localization characters with nontrivial inertia on their support.}
Fix an actual Hecke character $\psi$ satisfying
$\texttt{IsPsi}(\ell,Q)$: it has $\ell$-power order,
is trivial on $F$, is tame, ramifies exactly on $Q$,
and its entire local character is trivial above $\ell$.
Retain the prescribed separation of $Q$ from the type
and selected auxiliary primes and the common field of
values for $\mu_{0,j},\eta,\delta_\ell,\psi$.
GA8 constructs $\psi_\ell$ and a character $\chi_{\rm base}$ realizing
$B_r(\mu_{0,j}\delta_\ell\psi,\eta)$. With
$\varepsilon_\ell:G_E\to\mathbb Z_\ell^\times$
the actual cyclotomic character, define
\[
 \chi_c=\chi_{\rm base}\,\varepsilon_\ell^{-(c-1)}
       \quad(c\in\mathbb N),\qquad \chi_1=\chi_{\rm base}, \tag{GA.26}
\]
mapping $\varepsilon_\ell$ into $O^\times$. GA3 and
compatibility on $\mathbb Z_\ell$ prove continuity.
The integer exponent in (GA.26) also defines $c=0$;
later source uses may restrict to $1\leq c\leq t$.

GA7--GA10 show that $\psi_\ell$ has open kernel,
$\ell$-power order, anticyclotomic symmetry, exact
ramification set $Q$, and trivial action of every
decomposition group above $\ell$.
For each $P\in Q$, $P\nmid\ell$ follows already from
the local triviality of $\psi$ above $\ell$.
On its units the other factors of the block are trivial.
If $\chi_{\rm base}$ killed the whole inertia group, (GA.10) would
force $\psi_P$ to be trivial on units, contradicting
$P\in Q$. Hence $\chi_1$ is genuinely ramified at
every place in $Q$.

Away from $\ell$, the type primes, the selected auxiliary
prime and $Q$, put
$\vartheta=\mu_{0,j}\delta_\ell\psi\eta^{-1}$. The exact
Frobenius identity is
\[
 \iota\kappa(\chi_c(\sigma))
     =\vartheta(P)^{-1}N(P)^{-r/2-(c-1)}.               \tag{GA.27}
\]
Indeed (GA.11) and (GA.15) prove it for $c=1$.
Arithmetic Frobenius acts on every $\ell$-power root
of unity by its $N(P)$th power: reduction has that
action, and reduction is injective on these roots
when $P\nmid\ell$, since $X^{\ell^a}-1$ has unit
derivative at each root. This proves its cyclotomic
value is $N(P)$.
The source exponent at $c=1$ is
$((t-1)/2)+(1-n_V)/2=(t-n_V)/2=-r/2$,
since $n_V=r+t$. This proves every clause of
\texttt{LocalisationCharacterSpec}, conditional on
the actual $\psi,Q$; it does not construct
\texttt{IsPsi} for an arbitrary nonempty $Q$.

\paragraph{GA21: coefficient extension commutes with the normalized twist.}
Let $O_0,O$ have the same residue prime $\ell$, actual
readings $(\kappa_0,\iota_0)$ and $(\kappa,\iota)$,
and maps $f_0:K\to\operatorname{Frac}(O_0)$,
$f:K\to\operatorname{Frac}(O)$ reading the same
$\epsilon:K\to\mathbb C$. Suppose
$t:O_0\to O$ has the precise compatibility
\[
 y=f_0(x)\text{ in }\operatorname{Frac}(O_0)
       \quad\Longrightarrow\quad t(y)=f(x)
            \text{ in }\operatorname{Frac}(O)
 \quad(y\in O_0,\ x\in K).                             \tag{GA.28}
\]
Let $\chi_1,\chi,\psi_\ell$ be normalized realizations
of respectively $B_r(\mu_{0,j}\delta_\ell,\eta)$,
$B_r(\mu_{0,j}\delta_\ell\psi,\eta)$ and $\psi$.
Then
\[
                 \chi=t_*(\chi_1)\,\psi_\ell.           \tag{GA.29}
\]
Here $t$ is continuous by GA3 and both sides are continuous
characters into the same Hausdorff unit group $O^\times$.
For cofinally many places the inverse first-block value
is $\epsilon(x)$ for some $x\in K$, by GA12.
Equation (GA.11) identifies $\chi_1(\sigma)$ with
$f_0(x)$ in the fraction field, so (GA.28) identifies
its image under $t$ with $f(x)$.
Multiplying the two Frobenius values now gives the value
of $\chi$, since the two Hecke blocks differ by $\psi$.
RC4 proves (GA.29). This proves
\texttt{LocalisationTwistCharacterCompat} even when
the two algebraic complex readings differ away from $K$.
It does not assume that they commute on the entire
fraction field.

\paragraph{GA22: installed local characters, DeltaRead, and IsGalOf.}
The coefficient descent has a useful local form. Suppose
a supplied continuous $r_0:G_E\to O_0^\times$ has
cofinite normalized Frobenius readings in $K$ for a
specific algebraic Hecke character $\chi$. Let
$t:O_0\to O$ and the local field embedding
$f:K\to\operatorname{Frac}(O)$ satisfy the literal
integral-value compatibility (GA.28).
The target $\kappa,\iota$ must read $f$ as $\epsilon$.
Then $t_*r_0$ is the normalized realization of $\chi$.
To prove this, its cofinite readings put $\chi(P)$
cofinitely in $\epsilon(K)$; GA8 constructs a target
realization $r$. GA3 proves continuity of $t_*r_0$.
At the specified Frobenii its values and those of $r$
agree by the actual reading and coefficient equation;
conjugation transfers equality to every valuation above
the place. RC4 gives $r=t_*r_0$.

Apply this with the \emph{same} $D$, normalized
\texttt{SideB.chi.chi1} reading, local index, map
$T.\mathrm{twist}_2$, and $\operatorname{Sp.iq}$ of
the source \texttt{chi1\_loc\_isGal}. The retained
$\operatorname{AutSpec}$ identity
$Y.\mathrm{iotaQ}=\operatorname{Sp.iq}$ and its $LL$
reading give the target commuting square.
The equality $D.\mathrm{ntw}.\ell(T.\mathrm{twOf}(a))
=X.\mathrm{nd}.\ell(a)$ gives the same residue prime.
The three actual Hecke-character equalities of GA15
give exactly
$B_r(\mu_{0,j}\delta_{\ell(a)},\eta)$.
Thus this proves that literal endpoint, with all its
other premises preserved, although several are unused
by this character argument.

For any such local realization, GA7 immediately gives
\texttt{IsGalOf} on every set of places where its
specified Hecke factor is unramified and $P\nmid\ell$,
with the exponent in (GA.27). For a finite-order
installed $\Delta$ having the actual cofinite normalized
reading and \texttt{ValuesIn} in $LL_0$, the same
argument first identifies it with the realization of
$\delta_\ell$. At every unramified $P$, choose
$x_P\in LL_0$ reading $\delta_\ell(P)^{-1}$.
Injectivity gives the exact fraction-field equality for
$\Delta(\sigma)$, proving the full \texttt{DeltaRead}
clause, rather than only a cofinite version.
Transport to a local coefficient ring gives $\delta_\ell$
on the true side and $\delta_\ell^{-1}$ on the false
side: the supplied conjugate embedding conjugates its
root-of-unity value, and conjugation equals inversion.
Likewise an installed $\psi_\ell$ identified with the
same chosen $\psi$ gives the literal local
\texttt{NGalChrTied.psiL} equality. These implications
retain the equality of the chosen $\psi$ with the tuple
component; neither a tuple nor a local index is created
by them.

\paragraph{GA23: what can be installed, and the remaining simultaneous construction.}
For every already supplied twist index $i$, if the actual
$LL_0$ is star-stable and contains the HF values, GA16
constructs the entire source record
$\texttt{Triple.Model.Chars}(D,C,\operatorname{Art},i)$.
Choose a reading of its existing coefficient ring, as
in NC and RC8; star stability and the supplied equation
for the two coefficient embeddings give the conjugate
reading required in GA16. The resulting record has its
actual $\Delta$, both $\chi_1$ families, normalized
realization, open kernel, anticyclotomic law, exact
ramification and inertia order $15$, continuity and
the three inertia clauses (GA.21)--(GA.22).
GA14 removes dependence of these realization statements
on that initially chosen reading.

The literal \texttt{exists\_isTripleOf} can consequently
replace these character fields on a supplied twist
skeleton when $\mathrm{nd.Loc}$ is empty and the
specified tame array already satisfies
\texttt{TameTie}. It preserves the supplied set
$\mathrm{Tw}$, its coefficient rings, embeddings and
\texttt{NonzeroIndexPlace}. But its construction sets
all $\mathrm{Sig}$ and $\mathrm{abelian}_2$ to empty
and all form spaces to the zero module. The remaining
triple representation and occurrence assertions then
have empty domains. This is an accurately scoped
conditional endpoint, not a compatible arithmetic
realization of the form spaces used in RC or NC.

The unresolved construction is therefore precise:
one must supply nonempty twist and local index data on
the same $D$, cover every actual coefficient place by
the stipulated side embedding, install the normalized
character readings, and simultaneously satisfy the
form, Hecke, local and tuple identifications.
In particular the literal \texttt{SideB.cover} asserts
$\forall w\,\exists i$ with matching residue
characteristic and true-side valuation.
An empty $\mathrm{Tw}$ cannot satisfy it when a
coefficient place exists, and it cannot supply
$\mathrm{NChoice.twI}$.
GA13 constructs characters \emph{after} $i$ and its
reading have been supplied; it never reverses those
quantifiers. The local character constructions of
GA20--GA22 remove that part of the supplier once its
compatible inputs are present. Global reciprocity,
the nonempty indexing and field-containment/reading
conditions, the genuine representation suppliers and
the same-object $\mathrm{MtReal}/\mathrm{NExpTied}/
\mathrm{ColW}/\mathrm{Pre2}/\mathrm{St}$ joins remain
explicit. No conclusion of RH or GRH is asserted by
this section.
